% --- Core Mathematics Packages ---
\usepackage{amsmath, amssymb, amsthm}
\usepackage{mathtools} % For advanced formatting (e.g., \coloneqq for :=)

\usepackage{fancyhdr}
%\setlength{\headheight}{13.59999pt}
%\addtolength{\topmargin}{-1.59999pt}

% --- Typography and Page Layout ---
%\usepackage[margin=1in]{geometry}

\usepackage[a4paper, top=1in, bottom=1in, left=1.4in, right=1.4in]{geometry}

\usepackage{microtype} % Improves text justification and spacing
\usepackage{xcolor}    % Required for coloring links
\usepackage{hyperref}  % Makes the table of contents and cross-references clickable
\hypersetup{
    colorlinks=true,
    linkcolor=black,
    filecolor=magenta,      
    urlcolor=cyan,
    pdftitle={Mathematics Notes},
}
\usepackage{nameref} % Readable theorem referencing

\usepackage{subfiles}

\usepackage{framed}
\definecolor{shadecolor}{gray}{0.92}

\usepackage{tikz}
\usetikzlibrary{arrows.meta}
\usetikzlibrary{decorations.pathreplacing,calligraphy}


% --- Theorem Environments ---
% This links the numbering of theorems, lemmas, and definitions together.
\theoremstyle{definition}
\newtheorem{definition}{Definition}[chapter]
\newtheorem{example}[definition]{Example}

\theoremstyle{plain}
\newtheorem{theorem}[definition]{Theorem}
\newtheorem{lemma}[definition]{Lemma}
\newtheorem{corollary}[definition]{Corollary}
\newtheorem{proposition}[definition]{Proposition}

\theoremstyle{remark}
\newtheorem*{remark}{Remark} % The asterisk removes numbering for remarks


\newtheorem{exercise}{Exercise}[chapter]
\newtheorem*{exercise*}{Exercise}  % unnumbered, for one-off inline prompts


% --- Custom Shortcuts (Macros) ---
% Blackboard bold letters for standard sets
\newcommand{\R}{\mathbb{R}}
\newcommand{\N}{\mathbb{N}}
\newcommand{\Z}{\mathbb{Z}}
\newcommand{\Q}{\mathbb{Q}}
\newcommand{\C}{\mathbb{C}}

% Upright 'd' for differentials (Standard UK/Imperial convention)
\newcommand{\dd}{\mathrm{d}}